% Figure: aperture stop, entrance pupil, and the two defining rays. An off-axis
% object point sends a fan of rays into the system; the marginal ray grazes the
% edge of the aperture stop and fixes the cone the system accepts, while the chief
% ray passes through the centre of the stop and fixes the field. The image of the
% stop seen from object space is the entrance pupil.
\begin{figure}[htbp]
\centering
\begin{tikzpicture}[>=Stealth, scale=1.0]
% optical axis
\draw[enggray, dashed] (-0.4,0) -- (8.4,0);
% lens at x=3
\draw[very thick, engblue] (3,-1.5) -- (3,1.5);
\draw[engblue, ->] (3,1.5) -- (3,1.75);
\draw[engblue, ->] (3,-1.5) -- (3,-1.75);
\node[engblue, font=\scriptsize, anchor=south] at (3,1.75) {lens};
% aperture stop just behind lens at x=4.0 (two bars leaving a central gap)
\draw[very thick, engred] (4.0,0.6) -- (4.0,1.4);
\draw[very thick, engred] (4.0,-0.6) -- (4.0,-1.4);
\node[engred, font=\scriptsize, anchor=west] at (4.05,1.2) {aperture stop};
% object: off-axis arrow at x=0
\draw[engblack, ->, thick] (0,0) -- (0,1.0);
\node[font=\scriptsize, anchor=east] at (0,0.5) {object};
\filldraw[engblack] (0,1.0) circle [radius=1.0pt];
% marginal ray: from axial object point (0,0) grazing stop edge
\draw[engorange, thick] (0,0) -- (3,0.9);
\draw[engorange, thick] (3,0.9) -- (4.0,0.6);
\draw[engorange, ->, thick] (4.0,0.6) -- (6.6,-0.96);
\node[engorange, font=\scriptsize, anchor=south] at (1.5,0.5) {marginal ray};
% chief ray: from off-axis object tip through stop centre (4.0,0)
\draw[engamber, thick] (0,1.0) -- (3,0.55);
\draw[engamber, thick] (3,0.55) -- (4.0,0.0);
\draw[engamber, ->, thick] (4.0,0.0) -- (7.4,-1.87);
\node[engamber, font=\scriptsize, anchor=south east] at (2.6,0.9) {chief ray};
% image plane
\draw[engblue, thick] (7.4,-2.0) -- (7.4,0.4);
\node[engblue, font=\scriptsize, anchor=west] at (7.45,0.2) {image plane};
% entrance-pupil marker (virtual image of stop, drawn as dotted bars in object space)
\draw[engred, dotted, thick] (1.6,0.7) -- (1.6,1.4);
\draw[engred, dotted, thick] (1.6,-0.7) -- (1.6,-1.4);
\node[engred, font=\scriptsize, anchor=south] at (1.6,1.4) {entrance pupil};
\end{tikzpicture}
\caption[Aperture stop, entrance pupil, and the marginal and chief rays]{The
aperture stop is the physical opening that limits the cone of light the system
accepts from an axial point; the marginal ray, launched from the axial object
point, grazes the stop edge and fixes the numerical aperture and the depth of
field. The chief ray, launched from the edge of the field through the centre of
the stop, fixes the field of view and the vignetting. The image of the stop
formed by the optics ahead of it, seen from object space, is the entrance pupil,
the effective opening the object sees; the image of the stop formed by the optics
behind it is the exit pupil. Every aberration and the whole radiometry of the
system are referred to these two rays and these two pupils.}
\label{fig:vol04ch11:aperture-stop}
\end{figure}
